\section{Discussion}

\subsection{Root Cause Analysis}

The hypothesis failure has a clear mechanistic explanation: \textbf{feature saturation in pretrained models}.

CLIP ViT-B/16 was trained on 400 million image-text pairs. Both ``background'' and ``bird type'' are elementary visual categories well within CLIP's training distribution. Consequently:

\begin{enumerate}
    \item No emergence dynamics exist
    \item Probes converge immediately
    \item CV measures noise
\end{enumerate}

\subsection{Theoretical Interpretation}

Our negative result clarifies a fundamental distinction:

\begin{itemize}
    \item \textbf{Feature emergence}: A training-time phenomenon
    \item \textbf{Feature separability}: A representation property
\end{itemize}

C-sweep probing on frozen features measures separability, not emergence.

\subsection{Honest Limitations}

\begin{enumerate}
    \item \textbf{Single hypothesis tested}: Only H-E1 (existence) evaluated; intervention mechanism remains untested
    \item \textbf{Single dataset, single extractor}: Waterbirds with CLIP ViT-B/16
    \item \textbf{C-sweep as epoch proxy}: May not capture true learning dynamics
\end{enumerate}

\subsection{Implications for Future Work}

The negative result identifies necessary conditions for emergence-based detection:

\begin{enumerate}
    \item Training-time measurement
    \item Unfrozen representations
    \item Alternative signals (loss curves, gradient norms)
\end{enumerate}
